%Allgemeine Definitionen
\usepackage[ngerman]{babel}				% Eigenheiten der deutschen Sprache wie Umlaute und Silbentrennung einbinden (neue Rechtschreibung)
\usepackage[style=apa]{biblatex}					% quellenverzeichnisse (https://ctan.org/pkg/biblatex)
\usepackage{csquotes}					% Ergänzung für biblatex
\usepackage{fontspec}					%
\setmainfont{Arial}  					% Arial nur mit LUATEX
%\usepackage[scaled]{helvet}
\renewcommand{\familydefault}{\sfdefault}
\usepackage{xifthen}					%!Bedingte Compilierung ermgoeglichen (provides \isempty{};https://ctan.org/pkg/xifthen?lang=de)
\usepackage{scrhack}					%Warnungen die durch alte Pakete enststehen umgehen

%Formatierung
\usepackage{geometry}	
\usepackage{setspace}					% Vereinfachte Definition des Zeilenabstands
\usepackage{enumitem}					% Einstellung für Aufzählungen (https://www.namsu.de/Extra/pakete/Enumitem.html ,https://ctan.org/pkg/enumitem?lang=de)
\usepackage[table]{xcolor}				% Farbdefinitionen auch fuer Tabellen laden, muss vor pdfpages stehen (https://ctan.org/pkg/xcolor?lang=de)
\usepackage{caption}					%?Formatierungsmoeglichkeiten der Tabellen- und Bildunterschriften erweitern (https://ctan.org/pkg/caption?lang=de)
\usepackage{rotating}					% Rotieren von eigentlich allem (https://ctan.org/pkg/rotating?lang=de)
%\usepackage{background}					% Seitenhintergrund festlegen (muss nach xcolor geladen werden! https://ctan.org/pkg/background?lang=de)
\usepackage{listings}																					% Codedarstellung
%\usepackage{siunitx}					% Leichter Darstellung von Einheiten (https://ctan.org/pkg/siunitx)

% Tools
\usepackage{fixme}						% Kollaborationswerkzeug (toDoNotes) (https://ctan.org/pkg/fixme)
\usepackage{placeins}					% Verhindern das Tablen in falsche absätze rutschen (provides \FloatBarrier https://ctan.org/pkg/placeins?lang=de)

% Makros für Grafiken etc
\usepackage{graphicx}					% Einbinden von Bildern (https://ctan.org/pkg/graphicx?lang=de, https://www.namsu.de/Extra/pakete/Graphicx.html)
\usepackage{overpic}					% Übermalen von Bildern (https://www.namsu.de/Extra/pakete/Overpic.html, https://www.ctan.org/pkg/overpic)
\usepackage{tikz}						% Eigen Darstellungen "zeichen"
\usepackage{pdfpages}					% PDF Dateien (auch Teile) in das Dokument einfügen (https://www.namsu.de/Extra/pakete/Pdfpages.html, https://ctan.org/pkg/pdfpages?lang=de)

% Makros zu Tabellen
\usepackage{array}					% Tabellen ueber mehrere Seiten ermöglichen (https://www.ctan.org/pkg/longtable)
\usepackage{longtable}					% Tabellen ueber mehrere Seiten ermöglichen (https://www.ctan.org/pkg/longtable)
\usepackage[longtable]{multirow}		% Zeilen oder Spalten einer Tabelle zusammenfassen (https://www.namsu.de/Extra/pakete/Multirow.html, https://ctan.org/pkg/multirow?lang=de)
\usepackage{makecell}					% Mehrzeilige Zellen ermöglichen + anderes nettes (https://ctan.org/pkg/makecell?lang=de)
% Makros zum Beschriften und Referenzieren
\usepackage{hyperref}					% Hyperlinks verwenden, Sollte als letzes geladen werden! (https://www.namsu.de/Extra/pakete/Hyperref.html, https://ctan.org/pkg/hyperref?lang=de)
